\documentclass{article}
\usepackage[utf8]{inputenc}   % caracteres especiales (acentos, e�es)
\usepackage[spanish]{babel}     % varias definiciones para el espa�ol
\usepackage{graphicx}           % inserci�n de graficos
\usepackage[a4paper, left=2.54cm, right=2.54cm, top=2.54cm, bottom=2.54cm]{geometry}
\newenvironment{keywords}
  {\begin{quote}\small\textbf{Keywords:}}
  {\end{quote}}
\usepackage{lettrine}  % Paquete para la primera letra decorada
\begin{document}

\title{Trabajo Práctico Final: \\ 
	Organización de las Computadoras}
       

\author{
    Joaquín García A., \\
    Lucas Enedín García, \\
    Sebastián Ayrton Deharbe \\[1ex]
    \texttt{joa.garciq@gmal.com}, \\
    \texttt{lucasgarcia2297@gmail.com}, \\
    \texttt{sebasdeharbe@gmail.com} \\[1ex]
    \textit{UNL Facultad de Ingeniería y Ciencias Hídricas}
}

%\markboth{XX CONGRESO ARGENTINO DE LA FICH: COD-PAIS\_NRO}{}

\maketitle

\begin{abstract}
En este informe se presentará el diseño e implementacion de un procesador RISC - V monociclo visto en clase.
\end{abstract}

\begin{keywords}
Microarquitectura, RISC - V, procesador
\end{keywords}

\section{Introducción}

\lettrine{E}{n} este trabajo tenemos como objetivo construir un procesador que se basa en la arquitectura de conjunto de instrucciones RISC-V, un ISA libre y de código abierto. Siguiendo la filosofia RISC (conjunto de instrucciones reducido), la arquitectura lo que hace es implementar las instrucciones mas utilizadas para optimizar el rendimiento global.
Se construyó en Verilog un procesador RISC-V monociclo con la capacidad de ejecutar diversas instrucciones, entre ellas, sumas, restas, instrucciones de carga y almacenamiento de datos en memoria. El informe detalla las implementaciones necesarias para que este procesador pueda ejecutar el código de prueba en lenguaje .ASM que se nos ha proporcionado

\section{Desarrollo}
\lettrine{E}{l} sistema propuesto, consta dos módulos. El módulo de hardware, que es el procesador RV32I construido a partir de las guías prácticas 6 y 7 y de los diagramas y material teórico propuesto por la cátedra que a su vez esta dividido en Datapath y módulo memory. Además el sistema consta también de un módulo de software realizado en lenguaje ensamblador en RARS que incluye operaciones aritmético lógicas simples entre los valores inicializados en los registros.

\newpage

\subsection{Módulo de Hardware:}
El diseño y conexión de los módulos en un procesador monociclo de 32 bits es crucial para la ejecución eficiente de instrucciones. Cada uno de estos módulos tiene una función especifica que, en conjunto permite al procesador llevar a cabo operaciones lógicas, aritméticas y de control de flujo en un solo ciclo de reloj.

\begin{itemize}
\item \textbf{PC}: El contador de programa (PC) lleva la dirección de la proxima ejecución a ejecutar. Este módulo recibe el siguiente PC calculado, junto con la señal de reloj que determina cuando actualizarse
\item \textbf {Adder}: El sumador incrementa el valor del PC en 4, ya que cada instrucción ocupa 4 bytes en un sistema de 32 bits. Esto facilita el paso de una instrucción a la siguiente en memoria, pero también permite calcular la nueva dirección en el caso de una instrucción de salto.
\item \textbf {MUX1}: Este multiplexor selecciona entre dos posibles direcciones para el siguiente PC: una dirección secuencial (PC + 4) y una dirección de salto (calculada en base a la instrucción de salto).
\item \textbf{BR}: Este módulo es un arreglo de 32 registros de 32 bits. Almacena los datos temporales y permite la lectura y escritura durante la ejecución de instrucciones.
\item \textbf{SE}: Este módulo convierte un valor inmediato de menor tamaño en un valor de 32 bits, expandiéndolo según el valor de su bit más significativo para preservar el signo en operaciones de carga o salto. El extensor de signo toma un segmento de la instrucción y usa una señal de la unidad de control para generar un valor inmediato adecuado.
\item \textbf{ALU}: La ALU realiza operaciones aritméticas (como suma o resta) y lógicas (como AND, OR) utilizando dos entradas: una del banco de registros y otra seleccionada por MUX2. La ALU es controlada por la señal de ALUControl, que determina la operación específica que debe realizar, en función del tipo de instrucción.
\item \textbf{MUX2}: Selecciona entre dos posibles entradas para la ALU: Un valor leído del banco de registros o el valor inmediato generado por el extensor de signo. Esto permite flexibilidad para ejecutar instrucciones que usan un registro o un valor constante como entrada. 
\item \textbf{UC}: La unidad de control produce las señales de control que configuran los módulos en el datapath, tales como ALUControl para la ALU y las señales para la selección para los multiplexores. La UC determina la configuración adecuada para cada instrucción, a partir de los codigos f3, f7 y opcode de la instrucción
\item \textbf{mainDeco}: Se encarga de generar la señal de seleccion de operación para la ALU, en función del tipo de instrucción (códigos f3 y f7) y el opcode. Colabora con la unidad de control para establecer las configuraciones de operación especificas. 
\item \textbf{aluDeco}: Este modulo decodifica las señales de control y las traduce en instrucciones específicas que habilitan la ejecucion en cada componente del datapath. Cuando interpreta las señales provenientes de la UC, permite operaciones aritmeticas, lógicas o de control de flujo unicamente activando los módulos necesarios para la instrucción.
\item \textbf {MUX3: (4x1)}: Este módulo permite seleccionar entre múltiples fuentes de datos para determinar el valor final que será escrito o usado en el procesador. Tiene cuatro entradas y una salida, y usa señales de control para elegir cuál de las entradas debe pasar a la salida.
\end{itemize}

\newpage

\subsection{Módulo de memoria:}
Este módulo combina tanto la memoria de instrucciones  (IM) como la memoria de datos (DM) permitiendo al microprocesador almacenar y acceder a instrucciones yd atos necesarios para ejecutar un programa en un mismo bloque. Este módulo se encarga de gestionar la comunicacion entre el procesador y ambas memorias en un solo ciclo de reloj.
\begin{itemize}
\item \textbf{IM}: Este módulo se dedica a almacenar las instrucciones que el procesador va a ejecutar. Cada instruccion representa una operación que el procecesador debe realizar. En cada ciclo de reloj, el procesador accede a esta memoria para obtener la instruccion correspondiente a la direccion actual del contador de programa (PC).
\item \textbf{DM}: Esta memoria almacena los datos con los que trabaja el procesador durante la ejecucion del programa. Puede que contenga variables, resultados de operaciones o cualquier otro tipo de datos que la CPU necesite leer o escribir. Es independiente a la memoria de instrucciones y permite acceder a ambas en paralelo en un ciclo de reloj.
\end{itemize}

\begin{figure}[h]
  \centering
  \includegraphics[width=0.8\textwidth]{bloques_procesador.png}
  \caption{Diagrama de bloques del procesador RISCV32I.}
  \label{fig:mi_imagen}
\end{figure}

\newpage

\subsection{Módulo de Software}
El diseño y la interconexión de los módulos son esenciales para garantizar un funcionamiento correcto y eficiente del procesador, ya que cada módulo debe coordinarse con precisión en el flujo de datos y control. En nuestro caso, hemos desarrollado un procesador monociclo de 32 bits orientado a operaciones con datos enteros, donde cada ciclo de reloj permite procesar datos de 32 bits, lo que asegura precisión y rapidez. Para verificar la correcta implementación, usaremos un código de prueba proporcionado por la cátedra. Este código realiza operaciones aritméticas y lógicas básicas (como suma y resta) entre valores inicializados en el banco de registros, comprobando que los módulos como la ALU y el banco de registros funcionen correctamente. Luego, el procesador ejecuta dos bucles de prueba: un bucle \texttt{while}, que duplica un valor hasta alcanzar un límite, evaluando así la integridad del flujo en condiciones iterativas, y un bucle \texttt{for}, que realiza sumas repetidas. Finalmente, los resultados se almacenan en la memoria de datos y se recargan en los registros, comprobando la precisión en el almacenamiento y recuperación de datos, y asegurando que el procesador pueda completar la secuencia completa de operaciones.

\begin{figure}[h]
  \centering
  \includegraphics[width=0.8\textwidth]{Codigo_prueba.png}
  \caption{Codigo del programa proporcionado en ASM.}
  \label{fig:mi_imagen}
\end{figure}

\newpage

\section{Pruebas}
\lettrine{P}{a}ra validar el correcto funcionamiento del sistema se han implementado módulos de test bench que permitieron simular y verificar cada parte del diseño. Estas pruebas incluyeron tanto evaluacion de operaciones aritmético logicas para poder comprobar el funcionamiento del banco de registros y la ALU, así como la ejecución de instrucciones de salto y bucles. Los resultados obtenidos fueron observados y comparados con los valores esperados, confirmando la correcta interacción entre los módulos. Además estas pruebas nos han servido para comprender con mayor profundidad el funcionamiento interno de cada componente, proporcionando informacion valiosa sobre el rendimiento y la estabilidad del diseño. 
\begin{figure}[h]
  \centering
  \includegraphics[width=0.8\textwidth]{Prueba_completa.png}
  \caption{Test bench: Prueba completa.}
  \label{fig:mi_imagen}
\end{figure}


\section{Conclusiones}
\lettrine{E}{n} este trabajo se han aplicado conocimientos de arquitectura RISC-V para implementar un procesador monociclo, lo que nos permitió observar cómo una estructura modular, organizada y eficiente puede formar una unidad de procesamiento funcional. La experiencia nos brindó una compresión detallada del flujo de datos en la microarquitectura, asegurando la ejecución sincronizada de instrucciones en memoria. La naturaleza abierta de la arquitectura RISC también se destacó, mostrando su versatilidad para ser adaptada según las necesidades del usuario.

\begin{thebibliography}{9} 

    \bibitem{mano2003} M. Mano, \textit{Diseño digital}. Los Angeles: Pearson Education, Inc., 2003.

   \bibitem{patterson2018} Davit Patterson y Andrew Waterman, "Guia Práctica de RISC-V: El atlas de una arquitectura abierta.", 1era edicion, Straeberry Canyon LLC San Franciasco, Ca. USA 2018
    
    \bibitem{giovanini2024} L. Giovanini, ``Apuntes de cátedra,'' FICH UNL, Cátedra de Organización de las computadoras, 2022.
    
    \bibitem{padula2024} J. E. Padula, ``Apuntes de cátedra,'' FICH UNL, Cátedra de Organización de las computadoras, 2022.

\end{thebibliography}

\end{document}
